\subsection{Identifikasi Masalah}
\label{subsec:identifikasi-masalah}
Berdasarkan celah fungsional pada arsitektur DecMed yang telah diuraikan pada Subbab \ref{subsec:analisis-sistem-decmed}, kebutuhan operasional yang tampak pada Subbab \ref{subsec:analisis-skenario-layanan-klinis}, serta kewajiban perlindungan kerahasiaan dan pembatasan akses data pasien yang dibahas pada Subbab \ref{subsec:analisis-regulasi}, penelitian ini mengidentifikasi tiga masalah utama yang sekaligus menjadi fokus pembahasan pada Tugas Akhir ini. Analisis STRIDE pada Subbab \ref{subsec:analisis-stride-decmed} kemudian berperan sebagai penguat dari sisi keamanan, dengan menunjukkan bahwa ketiga masalah tersebut memang menimbulkan risiko yang nyata ketika sistem diterapkan.

\begin{enumerate}[leftmargin=*]
	\item Mekanisme kontrol akses pada DecMed belum mendukung otorisasi yang granular. Hak akses masih direpresentasikan secara statis dan belum dapat menyatakan batasan rinci seperti jenis data, aksi yang diperbolehkan, masa berlaku, maupun ruang lingkup delegasi. Kondisi ini belum selaras dengan kebutuhan layanan kesehatan yang menuntut penerapan prinsip \textit{need-to-know}.
	\item Skema penyimpanan data \textit{off-chain} pada DecMed belum dirancang secara tersegmentasi untuk mendukung penegakan kebijakan akses \textit{fine-grained}. Akibatnya, ketika akses diberikan kepada suatu objek data terenkripsi, sistem berpotensi memberikan cakupan akses yang lebih luas daripada kebutuhan aktual pengguna.
	\item Pendelegasian hak akses pada DecMed masih bergantung pada keterlibatan langsung pasien sebagai penerbit akses baru. Ketergantungan ini kurang sesuai dengan alur pelayanan kesehatan yang kolaboratif dan dinamis, karena delegasi antar tenaga medis seharusnya dapat berlangsung secara berantai dengan pembatasan yang makin ketat tanpa memerlukan intervensi ulang dari pasien.
\end{enumerate}

Ketiga masalah tersebut berkorespondensi langsung dengan rumusan masalah pada Bab \ref{chap:pendahuluan}, yaitu kebutuhan untuk merancang mekanisme kontrol akses granular berbasis Macaroons, menyusun skema penyimpanan \textit{off-chain} terenkripsi yang tersegmentasi, dan mendukung fitur atenuasi luring untuk delegasi akses berantai.
